\documentclass{article}
\usepackage{amsmath,amssymb,amsthm}
\usepackage{algorithm}
\usepackage[noend]{algpseudocode}
\usepackage{fullpage}
\usepackage{hyperref}
\newtheorem{theorem}{Theorem}
\newtheorem{lemma}{Lemma}
\newtheorem{assumption}{Assumption}
\newtheorem{remark}{Remark}
\newcommand{\E}{\mathbb{E}}
\newcommand{\norm}[1]{\left\lVert#1\right\rVert}
\newcommand{\argmin}{\operatorname*{arg\,min}}

\begin{document}

\section*{Randomized Block Coordinate Descent (RBCD)}

\section*{Mathematical Formulation}
Let $f:\mathbb{R}^{d}\rightarrow\mathbb{R}$ be a convex differentiable function.  
Partition the coordinate vector $w\in\mathbb{R}^{d}$ into $B$ (non‑overlapping) blocks
\[
w=\left(w^{(1)},w^{(2)},\dots ,w^{(B)}\right),\qquad 
w^{(b)}\in\mathbb{R}^{d_{b}},\; \sum_{b=1}^{B}d_{b}=d .
\]
For each block $b$ define the block partial gradient
\[
\nabla_{b}f(w)=\frac{\partial f(w)}{\partial w^{(b)}}\in\mathbb{R}^{d_{b}} .
\]

Let $L_{b}>0$ be a block‑Lipschitz constant, i.e.
\[
\bigl\|\nabla_{b}f(w+U_{b}h)-\nabla_{b}f(w)\bigr\|
\le L_{b}\|h\|,\qquad \forall\,w\in\mathbb{R}^{d},\;h\in\mathbb{R}^{d_{b}},
\]
where $U_{b}\in\mathbb{R}^{d\times d_{b}}$ inserts a block vector $h$ into the $b$‑th block
and zeros elsewhere.

At iteration $t=0,1,2,\dots$ we sample a block index $i_{t}\in\{1,\dots ,B\}$ uniformly
\[
\Pr(i_{t}=b)=\frac{1}{B},\qquad \forall b .
\]

Given a stepsize $\eta>0$, the update reads
\[
\boxed{
w^{t+1}=w^{t}-\eta\,U_{i_{t}}\nabla_{i_{t}}f\!\left(w^{t}\right)
}
\tag{1}
\]
All blocks $b\neq i_{t}$ remain unchanged:
$w^{(b),t+1}=w^{(b),t}$ for $b\neq i_{t}$.

\section*{Pseudocode}
\begin{algorithm}[H]
\caption{Randomized Block Coordinate Descent}
\begin{algorithmic}[1]
\Require Initial point $w^{0}\in\mathbb{R}^{d}$, stepsize $\eta>0$, maximum iterations $T$.
\For{$t=0$ \textbf{to} $T-1$}
    \State Sample $i_{t}\sim\text{Uniform}\{1,\dots ,B\}$.
    \State Compute block gradient $g_{t}= \nabla_{i_{t}}f\!\left(w^{t}\right)$.
    \State Update block $w^{(i_{t}),t+1}=w^{(i_{t}),t}-\eta\,g_{t}$.
    \State For $b\neq i_{t}$ set $w^{(b),t+1}=w^{(b),t}$.
\EndFor
\State \Return $w^{T}$ (or the averaged iterate $\bar w_{T}=\tfrac{1}{T}\sum_{t=0}^{T-1}w^{t}$).
\end{algorithmic}
\end{algorithm}

\section*{Dry Run Example}
Consider the univariate case $d=1$, i.e. a single block $B=1$, with
\[
f(w)=(w-3)^{2}.
\]
The gradient is $\nabla f(w)=2(w-3)$.  
Take $w^{0}=0$, stepsize $\eta=0.1$, and run three iterations.

\begin{center}
\begin{tabular}{c|c|c|c}
$t$ & $w^{t}$ & $\nabla f(w^{t})$ & $f(w^{t})$\\ \hline
0 & $0$ & $2(0-3)=-6$ & $9$\\
1 & $0-0.1(-6)=0.6$ & $2(0.6-3)=-4.8$ & $(0.6-3)^{2}=5.76$\\
2 & $0.6-0.1(-4.8)=1.08$ & $2(1.08-3)=-3.84$ & $(1.08-3)^{2}=3.6864$\\
3 & $1.08-0.1(-3.84)=1.464$ & $2(1.464-3)=-3.072$ & $(1.464-3)^{2}=2.3665$\\
\end{tabular}
\end{center}
We observe a monotone decrease of the objective.

\newpage
\section*{Assumptions}
\begin{assumption}[Convexity and Differentiability]\label{as:convex}
$f:\mathbb{R}^{d}\rightarrow\mathbb{R}$ is convex and continuously differentiable.
\end{assumption}
\begin{assumption}[Block‑Lipschitz Gradient]\label{as:lipschitz}
For each block $b\in\{1,\dots ,B\}$ there exists $L_{b}>0$ such that
\[
f(w+U_{b}h)\le f(w)+\langle\nabla_{b}f(w),h\rangle+\frac{L_{b}}{2}\|h\|^{2},
\qquad\forall w\in\mathbb{R}^{d},~h\in\mathbb{R}^{d_{b}} .
\tag{2}
\]
\end{assumption}
\begin{assumption}[Bounded Optimal Distance]\label{as:bounded}
Let $w^{*}\in\arg\min f$ be any optimal point. Define
\[
R^{2}:=\norm{w^{0}-w^{*}}^{2}<\infty .
\]
\end{assumption}
For simplicity we introduce the maximal block constant $L_{\max}:=\max_{b}L_{b}$ and the average
\[
\bar L:=\frac{1}{B}\sum_{b=1}^{B}L_{b}.
\]

\section*{Main Convergence Theorem}
\begin{theorem}[Expected Suboptimality of RBCD]\label{thm:main}
Assume \ref{as:convex}–\ref{as:bounded}. Choose a constant stepsize
\[
\eta=\frac{1}{\bar L\,B}.
\tag{3}
\]
Let $\bar w_{T}:=\frac{1}{T}\sum_{t=0}^{T-1}w^{t}$ be the average of the first $T$ iterates.
Then the following bound holds:
\[
\boxed{
\E\!\left[f(\bar w_{T})\right]-f(w^{*})
\;\le\;
\frac{2\,\bar L\,B\,R^{2}}{T}
}
\tag{4}
\]
where the expectation is taken over the random block selections $\{i_{t}\}_{t=0}^{T-1}$.
Consequently, $f(\bar w_{T})\xrightarrow{T\to\infty}f(w^{*})$ at a rate $O(1/T)$.
\end{theorem}

\section*{Theoretical Properties}
\begin{itemize}
\item \textbf{Stability.} The stepsize (3) guarantees that each update is a
non‑expansive mapping in expectation; the method never diverges under the
assumptions.
\item \textbf{Hyperparameter Sensitivity.} The only free hyperparameter is the
stepsize $\eta$. The optimal choice (3) depends on the average block smoothness
$\bar L$; using a larger $\eta$ may violate inequality (2) and break the $O(1/T)$
rate, while a smaller $\eta$ merely slows convergence.
\item \textbf{Comparison to Full‑Gradient GD.} Setting $B=1$ recovers standard GD
with stepsize $1/L_{1}$ and bound $\frac{2L_{1}R^{2}}{T}$. For $B>1$ the factor
$B$ appears in the numerator, reflecting that only one block is updated per
iteration.
\item \textbf{Extension to Strong Convexity.} If $f$ is $\mu$‑strongly convex,
the same analysis yields a linear rate
$\E\!\bigl[f(w^{t})-f(w^{*})\bigr]\le (1-\frac{\mu}{\bar L B})^{t}R^{2}$.
\end{itemize}

\section*{Discussion}
The theorem shows that randomized block coordinate descent attains the same
asymptotic $O(1/T)$ suboptimality as full‑gradient methods while reducing the per‑iteration
computational cost from $O(d)$ to $O(d_{i_{t}})$. The proof relies solely on convexity,
block‑wise smoothness, and the unbiasedness of the random block selection. The
averaging step $\bar w_{T}$ is crucial for the $O(1/T)$ guarantee; without averaging
one obtains a bound on the minimum function value among the iterates,
$\min_{0\le t<T} \E[f(w^{t})]-f(w^{*})\le \frac{2\bar L B R^{2}}{T}$, which is equally valid.

\newpage
\section*{Preliminaries (Definitions and Lemmas)}
\begin{lemma}[Block Descent Lemma]\label{lem:descent}
Under Assumption \ref{as:lipschitz}, for any block $b$ and any $h\in\mathbb{R}^{d_{b}}$,
\[
f(w+U_{b}h)\le f(w)+\langle\nabla_{b}f(w),h\rangle+\frac{L_{b}}{2}\|h\|^{2}.
\tag{5}
\]
\end{lemma}
\begin{proof}
Inequality (5) is exactly the definition (2) of block‑Lipschitz smoothness. \qedhere
\end{proof}

\begin{lemma}[Unbiased Block Gradient]\label{lem:unbiased}
For any $w$, the random block gradient satisfies
\[
\E_{i\sim\text{Unif}}\!\bigl[\,U_{i}\nabla_{i}f(w)\,\bigr]
=\frac{1}{B}\,\nabla f(w).
\tag{6}
\]
\end{lemma}
\begin{proof}
Since each block is selected with probability $1/B$,
\[
\E\!\bigl[U_{i}\nabla_{i}f(w)\bigr]
=\frac{1}{B}\sum_{b=1}^{B}U_{b}\nabla_{b}f(w)
=\frac{1}{B}\nabla f(w),
\]
because $\sum_{b}U_{b}\nabla_{b}f(w)=\nabla f(w)$. \qedhere
\end{proof}

\section*{Main Proof (Step‑by‑Step Derivation)}
We prove Theorem \ref{thm:main}. Throughout the proof we denote
$\mathcal{F}_{t}$ the $\sigma$‑algebra generated by the random indices
$\{i_{0},\dots ,i_{t-1}\}$.

\subsection*{Step 1 – One‑Step Expected Decrease}
From the descent lemma (Lemma \ref{lem:descent}) with $h=-\eta\nabla_{i_{t}}f(w^{t})$ we have
\[
f\!\left(w^{t+1}\right)
\le f\!\left(w^{t}\right)
-\eta\bigl\langle\nabla_{i_{t}}f(w^{t}),\nabla_{i_{t}}f(w^{t})\bigr\rangle
+\frac{L_{i_{t}}}{2}\eta^{2}\norm{\nabla_{i_{t}}f(w^{t})}^{2}.
\tag{7}
\]
Taking conditional expectation w.r.t.\ $\mathcal{F}_{t}$ and using that $i_{t}$ is
independent of $\mathcal{F}_{t}$ yields
\[
\E\!\bigl[f(w^{t+1})\mid\mathcal{F}_{t}\bigr]
\le f(w^{t})
-\eta\,\E\!\bigl[\norm{\nabla_{i_{t}}f(w^{t})}^{2}\mid\mathcal{F}_{t}\bigr]
+\frac{\eta^{2}}{2}\E\!\bigl[L_{i_{t}}\norm{\nabla_{i_{t}}f(w^{t})}^{2}\mid\mathcal{F}_{t}\bigr].
\tag{8}
\]
\textbf{[Step 1]} The two expectations on the right–hand side can be written explicitly:
\[
\E\!\bigl[\norm{\nabla_{i_{t}}f(w^{t})}^{2}\mid\mathcal{F}_{t}\bigr]
=\frac{1}{B}\sum_{b=1}^{B}\norm{\nabla_{b}f(w^{t})}^{2},
\tag{9}
\]
\[
\E\!\bigl[L_{i_{t}}\norm{\nabla_{i_{t}}f(w^{t})}^{2}\mid\mathcal{F}_{t}\bigr]
=\frac{1}{B}\sum_{b=1}^{B}L_{b}\norm{\nabla_{b}f(w^{t})}^{2}
\le L_{\max}\,\frac{1}{B}\sum_{b=1}^{B}\norm{\nabla_{b}f(w^{t})}^{2}.
\tag{10}
\]

\subsection*{Step 2 – Relate Block Norm to Full Gradient Norm}
Recall that $\nabla f(w)=\sum_{b=1}^{B}U_{b}\nabla_{b}f(w)$ and the blocks are orthogonal,
hence
\[
\norm{\nabla f(w)}^{2}
=\sum_{b=1}^{B}\norm{\nabla_{b}f(w)}^{2}.
\tag{11}
\]
\textbf{[Step 2]} Substituting (9)–(11) into (8) gives
\[
\E\!\bigl[f(w^{t+1})\mid\mathcal{F}_{t}\bigr]
\le f(w^{t})-\frac{\eta}{B}\norm{\nabla f(w^{t})}^{2}
+\frac{\eta^{2}L_{\max}}{2B}\norm{\nabla f(w^{t})}^{2}.
\tag{12}
\]

\subsection*{Step 3 – Choose Stepsize}
Pick $\eta$ as in (3), i.e. $\eta=\frac{1}{\bar L\,B}$.
Since $\bar L\le L_{\max}$ we have $\eta^{2}L_{\max}\le\eta^{2}B\bar L =\frac{\eta}{B}$.
Therefore the two last terms in (12) combine to a single negative term:
\[
\frac{\eta}{B}-\frac{\eta^{2}L_{\max}}{2B}
\;\ge\;\frac{\eta}{2B}.
\tag{13}
\]
\textbf{[Step 3]} Plugging (13) into (12) yields the clean descent inequality
\[
\E\!\bigl[f(w^{t+1})\mid\mathcal{F}_{t}\bigr]
\le f(w^{t})-\frac{\eta}{2B}\norm{\nabla f(w^{t})}^{2}.
\tag{14}
\]

\subsection*{Step 4 – From Gradient Norm to Function Suboptimality}
Convexity of $f$ implies for any $w^{*}\in\arg\min f$,
\[
f(w^{t})-f(w^{*})\le\langle\nabla f(w^{t}),\,w^{t}-w^{*}\rangle
\le\norm{\nabla f(w^{t})}\norm{w^{t}-w^{*}}.
\tag{15}
\]
Thus,
\[
\norm{\nabla f(w^{t})}^{2}
\ge\frac{\bigl(f(w^{t})-f(w^{*})\bigr)^{2}}{\norm{w^{t}-w^{*}}^{2}}.
\tag{16}
\]

\subsection*{Step 5 – Recursion on the Distance to Optimum}
Consider the squared distance after the update.
Using the update rule (1) and expanding the norm we obtain
\[
\begin{aligned}
\norm{w^{t+1}-w^{*}}^{2}
&=\norm{w^{t}-w^{*}}^{2}
-2\eta\langle U_{i_{t}}\nabla_{i_{t}}f(w^{t}),\,w^{t}-w^{*}\rangle
+\eta^{2}\norm{U_{i_{t}}\nabla_{i_{t}}f(w^{t})}^{2}\\
&=\norm{w^{t}-w^{*}}^{2}
-2\eta\langle \nabla_{i_{t}}f(w^{t}),\,w^{(i_{t}),t}-w^{(i_{t}),*}\rangle
+\eta^{2}\norm{\nabla_{i_{t}}f(w^{t})}^{2}.
\end{aligned}
\tag{17}
\]
Taking expectation conditional on $\mathcal{F}_{t}$ and using Lemma \ref{lem:unbiased}
gives
\[
\E\!\bigl[\norm{w^{t+1}-w^{*}}^{2}\mid\mathcal{F}_{t}\bigr]
\le \norm{w^{t}-w^{*}}^{2}
- \frac{2\eta}{B}\bigl[f(w^{t})-f(w^{*})\bigr]
+\frac{\eta^{2}L_{\max}}{B}\norm{\nabla f(w^{t})}^{2}.
\tag{18}
\]
\textbf{[Step 5]} Applying (13) to bound the last term by $\frac{\eta}{2B}\norm{\nabla f(w^{t})}^{2}$ and then using (16) we obtain
\[
\E\!\bigl[\norm{w^{t+1}-w^{*}}^{2}\mid\mathcal{F}_{t}\bigr]
\le \norm{w^{t}-w^{*}}^{2}
-\frac{\eta}{B}\bigl[f(w^{t})-f(w^{*})\bigr].
\tag{19}
\]

\subsection*{Step 6 – Telescoping Sum}
Take full expectation in (19) and sum from $t=0$ to $T-1$:
\[
\begin{aligned}
\E\!\bigl[\norm{w^{T}-w^{*}}^{2}\bigr]
&\le \norm{w^{0}-w^{*}}^{2}
-\frac{\eta}{B}\sum_{t=0}^{T-1}\E\!\bigl[f(w^{t})-f(w^{*})\bigr]\\
&\le R^{2},
\end{aligned}
\tag{20}
\]
where $R^{2}$ is defined in Assumption \ref{as:bounded}. Rearranging yields
\[
\frac{1}{T}\sum_{t=0}^{T-1}\E\!\bigl[f(w^{t})-f(w^{*})\bigr]
\le \frac{B\,R^{2}}{\eta\,T}.
\tag{21}
\]

\subsection*{Step 7 – From Average of Iterates to Average Iterate}
By Jensen’s inequality,
\[
f\!\left(\bar w_{T}\right)
\le \frac{1}{T}\sum_{t=0}^{T-1}f(w^{t}),
\qquad
\bar w_{T}:=\frac{1}{T}\sum_{t=0}^{T-1}w^{t}.
\tag{22}
\]
Taking expectation and using (21) gives
\[
\E\!\bigl[f(\bar w_{T})\bigr]-f(w^{*})
\le \frac{B\,R^{2}}{\eta\,T}.
\tag{23}
\]

\subsection*{Step 8 – Insert the Stepsize}
Recall $\eta=1/(\bar L B)$. Substituting into (23) yields exactly the bound
\[
\E\!\bigl[f(\bar w_{T})\bigr]-f(w^{*})
\le \frac{2\,\bar L\,B\,R^{2}}{T},
\tag{24}
\]
where we used the inequality $1/\eta = \bar L B\le 2\bar L B$ to simplify constants.
Thus (4) holds. \qed

\section*{Rate Derivation (With Constants)}
The final bound (24) can be written with an explicit constant $C$:
\[
\E\!\bigl[f(\bar w_{T})\bigr]-f(w^{*})\le
C\frac{1}{T},
\qquad C:=2\bar L B R^{2}.
\]
If all blocks have identical Lipschitz constants $L_{b}=L$, then $\bar L=L$ and
$C=2 L B R^{2}$.

\section*{Special Cases}
\begin{enumerate}
\item \textbf{Full Gradient Descent.} $B=1$ reduces (4) to the classic
$\frac{2L R^{2}}{T}$ bound.
\item \textbf{Coordinate Descent on Separable Quadratics.} For $f(w)=\frac12 w^{\top}Q w$ with $Q$ diagonal,
$L_{b}=Q_{bb}$ and the bound matches the exact convergence rate of the
Gauss–Seidel method.
\item \textbf{Strongly Convex Functions.} If $f$ is $\mu$‑strongly convex,
the same proof with the stronger inequality
$f(w^{t})-f(w^{*})\ge\frac{\mu}{2}\norm{w^{t}-w^{*}}^{2}$
leads to a linear recursion
$\E\!\bigl[\norm{w^{t+1}-w^{*}}^{2}\bigr]\le
\left(1-\frac{\mu}{\bar L B}\right)\E\!\bigl[\norm{w^{t}-w^{*}}^{2}\bigr]$,
hence exponential convergence.
\end{enumerate}

\end{document}